% Abhakseite „Das kann ich“ – Zone nur im Gesamt (\mitzone)
%% TODO Ich-kann-Sätze: Platzhalter wie in den Aufgabendateien
\begin{abhakseite}
\ifdefined\mitzone
\abhakgruppe{Kennst du schon}
\abhak{1}{Sinus, Kosinus und Tangens als Seitenverhältnisse im rechtwinkligen Dreieck, Wert zu einem Winkel mit dem Taschenrechner im Grad-Modus}
\abhak{2}{Taschenrechner}
\abhak{3}{Koordinatensystem mit vier Quadranten, Punkte eintragen und ablesen, Achsen benennen und einteilen}
\abhak{4}{Kreis mit Mittelpunkt und Radius zeichnen, Umfang mit Pi berechnen, Mittelpunktswinkel und Kreisbogen als Anteil des Vollkreises}
\abhak{5}{Winkel messen, zeichnen und benennen, Vollwinkel, gestreckter und rechter Winkel, Winkel über neunzig Grad als stumpfe Winkel}
\abhak{6}{Bruchteile eines Vollkreises als Bruch und als Dezimalzahl (halb, viertel, drittel)}
\abhak{7}{Wertetabelle anlegen, Wertepaare als Punkte eintragen, den Graphen als durchgehende Linie zeichnen, Funktionswert und Argument ablesen, Punktprobe}
\abhak{8}{Achseneinteilung wählen, wenn beide Achsen verschiedene Einheiten und Schrittweiten haben}
\abhak{9}{Achsensymmetrisch und punktsymmetrisch als Wörter, Symmetrieachse und Symmetriezentrum am Graphen}
\abhak{10}{Parameter in einer Funktionsgleichung erkennen und ihre Wirkung auf den Graphen beschreiben (Streckung, Stauchung, Verschiebung), Graph zu Gleichung zuordnen}
\abhak{11}{Größtwert, Kleinstwert und Spannweite einer Tabelle entnehmen, Mittelwert bilden}
\abhak{12}{Taschenrechner – Fehler finden}
\abhak{13}{Taschenrechner – selbst rechnen}
\fi
\abhakgruppe{Einheit 1 · Einheitskreis und Bogenmaß}
\abhak{14}{Einheitskreis und Bogenmaß}
\abhak{15}{Einheitskreis und Bogenmaß – Fehler finden, Begründen}
\abhakgruppe{Einheit 2 · Sinus- und Kosinusfunktion und ihre Merkmale}
\abhak{16}{Sinus- und Kosinusfunktion}
\abhak{17}{Sinus- und Kosinusfunktion – Fehler finden, Begründen, Darstellungswechsel}
\abhakgruppe{Einheit 3 · Parameter und Transformationen}
\abhak{18}{Parameter und Transformationen}
\abhak{19}{zwei Graphen mit verschiedenen Parametern vergleichen}
\abhak{20}{Parameter und Transformationen – Fehler finden, Begründen}
\abhakgruppe{Einheit 4 · Periodische Vorgänge modellieren}
\abhak{21}{Periodische Vorgänge modellieren}
\abhak{22}{Periodische Vorgänge modellieren – Fehler finden, Begründen}
\end{abhakseite}
